\documentclass{article}
\usepackage{pathfinder-readable}
\date{}
\begin{document}
\section{Certifying polarisation observables before transferring DCQAOA}

Julien, the useful connection between EMLE and digitised-counterdiabatic QAOA is a testable question: can a quantum optimiser improve a specified polarisation observable after accounting for binary encoding and comparing against suitable classical methods? The approved research supports conditional certificates and synthetic counterexamples. Production replacement of the classical dipole solve should remain paused because the publication-specific energy convention, matched molecular inputs and useful accuracy-versus-cost regime remain unresolved. The next step is a publication-matched response fixture, followed by an observable-level benchmark.

The account distinguishes results inherited from the papers and branch investigations, connections established in the frozen joint investigation, and prospective benefits that remain conjectural. References labelled S below identify the approved generated research supplement, with stable member paths and ledger references in the appendix. S is reviewable campaign evidence, not independent molecular corroboration.

\section{Physical problem and proposed quantum method}

Morado and colleagues' \emph{Enhancing Electrostatic Embedding for ML/MM Free Energy Calculations}, Q, studies electrostatic machine learning embedding for hydration free energies, including calibration of static electrostatics and induction. Ding, Huang and Yuan's \emph{Molecular docking via quantum approximate optimization algorithm}, P, studies QAOA variants on pharmacophore contact graphs. Their work identities are \url{https://doi.org/10.1021/acs.jctc.5c01464} and \url{https://doi.org/10.1103/PhysRevApplied.21.034036}.

On Q's training set, patched EMLE gives a hydration RMSE of 1.63 kcal/mol, compared with 0.80 for OpenFF/TIP3P. Empirical attenuation with inverse dielectric 0.6 reduces the former to 0.74 kcal/mol. Q reports the scaled predictions as statistically indistinguishable from OpenFF/TIP3P and mechanical-embedding ML/MM at the 95\% confidence level. These inherited results concern physical model calibration, not quantum optimisation. Q hypothesises double-counting of induction in the combined model as the main cause of excessive hydration. Its discussion of a training bottleneck concerns reference-data generation; it does not establish the dipole solve as the bottleneck. Better numerical minimisation alone is therefore not an evidenced remedy for the model discrepancy. See Q, manuscript pp. 17--18 and 24--30.

For a fixed configuration, let $N$ be the number of polarisable sites represented. Each site's induced dipole has three Cartesian components, giving a response matrix $A$ of dimension $d=3N$. Q's induced-dipole system, manuscript p. 13, Eq. 11, connects this matrix to molecular quantities: in the isotropic model, its diagonal three-by-three self-response blocks are $\alpha_j^{-1}I_3$, where $\alpha_j$ is the atomic polarisability and $I_3$ is the identity matrix. Its off-diagonal blocks are minus the Thole-damped dipolar interaction tensors $T_{jk}$. The field $E$ collects the electric fields generated by the MM charges at the sites.

The conditional starting point assumes a real symmetric positive-definite response matrix and defines
\begin{equation}
U(\mu;E)=\frac12\mu^TA\mu-E^T\mu,\qquad A\mu_*=E,\qquad U_*=-\frac12E^TA^{-1}E.
\end{equation}
The molecular identification of the blocks does not validate a publication-matched matrix. Symmetry, positive definiteness, the field-scaling convention and compatibility with Q's historical implementation must still be verified on the matched fixture. The proposed encoding is $\mu=\ell+Bx$, with binary components of $x$. With $b$ bits per dipole component, the direct representation uses $3Nb$ binary variables and qubits. Substitution produces a quadratic binary objective with potentially dense couplings. Even diagonal $A$ couples different bits encoding the same component through its self-energy. See S, branch-1 energy and operator notes.

P's simulations use independent gate parameters. Its layers contain ZZ and Z cost rotations, X mixer rotations and, for DCQAOA, an additional local Y layer. Classical optimisation adjusts the parameters; computational-basis measurements supply bitstrings. P reports selected 6-, 8- and 12-qubit docking simulations, not hydration calculations or demonstrated hardware advantage. The inherited ansatz audit shows that unrestricted independent local rotations, with ZZ angles set to zero, allow both circuit families to prepare every basis-state optimum at one layer. This establishes representability, not efficient discovery. See P, Section II.C, Section III and Figure 9; S, branch-1 product-control note and branch-2 ansatz results.

\section{Separate the objective, decoded output and observable}

The exact proposed snapshot target is conditional on the implementation's energy contract:
\begin{equation}
A\mu_*=E_s,\qquad H_*=-\frac12E_0^T\mu_*.
\end{equation}
Here $E_s$ is the solve field and $E_0$ is the energy-contraction field. An inspection reported in S found distinct damped solve and bare readout fields in a later EMLE-engine revision. That revision has not been established as Q's historical implementation. A matched fixture must determine whether the fields coincide, are uniformly proportional, or require a genuinely mixed-field treatment.

\paragraph{Equal fields.}
If $E_s=E_0=E$, the conventional encoded quantum objective is the expectation of the diagonal cost $U(\ell+Bx;E)$. A decoded selected bitstring supplies a vector $\tilde\mu$; the compatible variational readout is $U(\tilde\mu;E)$, compared with $H_*=U_*$. The expected cost used in training is not automatically the energy of that selected vector or of the decoded mean. For $r=A\tilde\mu-E$, the inherited audit gives
\begin{align}
U(\tilde\mu;E)-U_*&=\frac12r^TA^{-1}r\geq0,\\
U(\tilde\mu;E)+\frac12E^T\tilde\mu&=\frac12\tilde\mu^Tr.
\end{align}
The raw half-interaction readout $-E^T\tilde\mu/2$ instead has signed, first-order error in dipole error. It should be recorded separately when testing a proposed replacement interface. It does not inherit the variational certificate.

\paragraph{Verified uniform scaling.}
If $E_s=cE_0$ with constant $c>0$, the compatible functional is
\begin{equation}
W_c(\mu)=\frac{\mu^TA\mu}{2c}-E_0^T\mu.
\end{equation}
Its minimiser is $cA^{-1}E_0$ and its minimum is $H_*$. The proposed quantum training cost is the encoded expectation of $W_c$; a selected decoded vector is reported through $W_c$, with the raw half-interaction retained as a separate diagnostic. The alternative solve objective $\mu^TA\mu/2-cE_0^T\mu$ has the same minimiser but a minimum multiplied by $c$. Matching dipoles therefore does not establish matching energies. These equal-field and scaling contracts are inherited conditional results in S, branch-1 derivations and branch-2 implementation audit.

\subsection{Joint contribution: mixed-field readout}

For genuinely different fields, the joint investigation introduces an approximate adjoint $z$ for $Az_*=E_0$. Define $r_s=A\mu-E_s$ and $r_0=Az-E_0$. Its corrected readout is
\begin{align}
L(\mu,z)&=-\frac12E_0^T\mu+\frac12z^Tr_s\\
&=\frac12\left(z^TA\mu-E_0^T\mu-E_s^Tz\right),\\
L-H_*&=\frac12r_0^TA^{-1}r_s,\\
|L-H_*|&\leq\frac{\Vert r_0\Vert\Vert r_s\Vert}{2a},\qquad 0<a\leq\lambda_{\min}(A).
\end{align}
The norm is Euclidean and $a$ must be a certified positive lower spectral bound. The correction is exact if either vector is exact. Equal fields with $z=\mu$ recover $U$; uniform scaling with $z=\mu/c$ recovers $W_c$.

For an operational benchmark, obtain the primal approximation by solving the SPD problem with field $E_s$, and the adjoint approximation by solving the SPD problem with field $E_0$. If a solve uses the proposed quantum encoding, its conventional training objective is the expectation of the corresponding quadratic $U(\mu;E_s)$ or $U(z;E_0)$. Decode the resulting vectors under a declared output rule, then evaluate both the raw readout $-E_0^T\mu/2$ and the corrected readout $L(\mu,z)$ against the same reference $H_*$. Neither individual quadratic minimum is the general mixed-field target.

This specifies the logical separation of solving and reporting, not a completed quantum implementation. Whether the adjoint is classical or quantum, its accuracy, encoding, allocation of shots and reuse across evaluations remain choices to be fixed before a trial. The approved research does not supply a trained two-solve quantum workflow. A classical exact adjoint is a useful diagnostic but its cost belongs in the comparison.

In particular, $L$ is evaluated after obtaining approximate solutions; it is not jointly minimised. Its unrestricted Hessian is
\begin{equation}
\nabla^2L=\begin{pmatrix}0&A/2\\A/2&0\end{pmatrix},
\end{equation}
with eigenvalues $\pm\lambda_i(A)/2$. The stationary point is a saddle. In the retained scalar test with $A=1$, $E_s=1$, $E_0=2$ and both variables on the integer grid from $-3$ to $3$, joint minimisation gives $-6$ at $(\mu,z)=(3,-3)$, whereas the intended stationary value is $-1$ at $(1,2)$.

The frozen exact-rational mixed-field tests verify the identity, bound, scaling restrictions and both error signs. Correction can worsen a particular absolute error: at spacing $1/2$, raw error is $-0.0010479$ and corrected error is $+0.0014471$. The prescribed fields are not shown to arise from a shared MM environment. These are dimensionless arithmetic checks, not molecular evidence. S attributes adjoint residual correction to established work by Venditti and Darmofal, Sections 2.1 and 3.1, \url{https://acdl.mit.edu/darmofalpubs/JCP_vol187_p22.pdf}. The joint contribution is its use as a compatibility test for this transfer, not invention of the method.

\section{Encoding barriers and classical controls}

The inherited branch calculations show that well-conditioned continuous response problems can acquire many strict single-bit minima after binary encoding. The joint investigation connects those traps to product-state training. If $x^{(j)}$ is obtained by flipping bit $j$, small independent rotations around a basis string give
\begin{equation}
C(\delta)=f(x)+\frac14\sum_j\left[f(x^{(j)})-f(x)\right]\delta_j^2+O(\Vert\delta\Vert^4).
\end{equation}
The angle Hessian is diagonal with entries equal to half the single-bit energy gaps. A strict bit minimum is therefore a strict local minimum in this product-angle family. This is not a Hessian theorem for the full entangling P circuit.

The dimensionless example in S uses
\begin{equation}
U(\mu)=\frac12\mu^2+0.51\mu,\qquad \mu=-2+x_0+2x_1.
\end{equation}
The continuous optimum is $-0.51$, with energy $-0.13005$. The grid value zero is a strict bit minimum although the adjacent dipole level $-1$, with energy $-0.01$, is better. Along independent probabilities $p_0=t$, $p_1=1-t$, the mean is $-t$ and
\begin{equation}
C_{\mathrm{product}}(t)=1.99t-2t^2.
\end{equation}
This prescribed path reaches a barrier of $39601/80000=0.4950125$ at $t=199/400$. It is not a minimum-barrier result over all paths. A correlated classical mixture supported on zero and minus one has the same mean but expectation $-0.01t$, which decreases monotonically. Correlations can remove this demonstrated encoding barrier classically. No retained experiment shows that a trained P circuit discovers the path at favourable cost.

The inherited two-site controls impose explicit undamped isotropic dipole blocks. The axis-aligned system has condition number $127/123$ and 64, 64, 729 and 729 strict bit minima at spacings $1,1/2,1/4,1/8$, with one strict adjacent-level minimum at each spacing. Full enumeration checks the 12- and 18-bit cases; finer results use exact Cartesian-pair factorisation. A rotated system has condition number approximately 1.01207 and 64 strict bit minima at each spacing, with full enumeration but trimmed finer-grid domains. The differing counts concern differing geometries and encodings, so they do not conflict or establish a common asymptotic scaling law. Both sequences retain an energy-error plateau. A separate field on the axis-aligned matrix gives rounding excess $0.006968$, so solve-and-round is not generally an exact discrete solver. See S, branch-3 two-site results and ledger entries 21--27.

The joint continuation retains 64 strict bit/product minima and one strict adjacent-level minimum at coupling strengths $0,1/4,1/2,1$. Its field changes to keep the continuous solution fixed. This is not a fixed-field molecular experiment.

These results motivate component-level classical controls alongside independent-bit optimisation. For independent component distributions with means $m_i$ and second moments $q_i$,
\begin{equation}
\mathbb E U=U(m)+\frac12\sum_i A_{ii}(q_i-m_i^2).
\end{equation}
Bits within a component may be correlated. Holding other components fixed gives the conditional energy
\begin{equation}
\frac12 A_{ii}v^2+v\left(\sum_{j\neq i}A_{ij}m_j-E_i\right).
\end{equation}
On a uniform bounded grid, its best level is the clipped nearest level to
\begin{equation}
\frac{E_i-\sum_{j\neq i}A_{ij}m_j}{A_{ii}}.
\end{equation}
Component updates can cross binary carries directly; global discrete convergence is not guaranteed. The comparison should retain a continuous direct or suitable iterative solve, force-consistent classical approximations, solve-and-round, exact small-grid minima, analytic product optimisation and these component updates. See S, joint investigation.

\subsection{What a counterdiabatic interpretation requires}

With independently declared conventions
\begin{align}
H_C&=c_0I+\sum_i h_iZ_i+\sum_{i<j}J_{ij}Z_iZ_j,\\
H_M&=-\sum_i g_iX_i,
\end{align}
and $H(\lambda)=(1-\lambda)H_M+\lambda H_C$, the inherited operator audit gives
\begin{align}
i[H,\partial_\lambda H]&=-2\sum_i g_ih_iY_i\\
&\quad-2\sum_{i<j}J_{ij}(g_iY_iZ_j+g_jZ_iY_j).
\end{align}
P's added local Y layer is not the full first-commutator direction for a generic encoded response. It remains a variational extension, and composition of the circuit can produce nonlocal effects. A driving interpretation requires an interpolation, schedule, coefficient prescription and discretisation control. Attributing optimisation gains to transition suppression additionally needs dynamical evidence.

For an operator $O$ on $L$ qubits, the normalised Hilbert--Schmidt norm is $\|O\|_{\rm HS}^2=2^{-L}\operatorname{Tr}(O^\dagger O)$. Here the pair part of the first-commutator operator is partitioned into Pauli terms whose two bits encode the same dipole component and terms whose bits encode different components. These orthogonal Pauli terms contribute their squared coefficients to the norm. With unit mixer strengths, each unordered bit pair contributes $8J_{ij}^2$. In the joint continuation, summing separately over these two classes gives squared norms 12 within components and at most 0.0048 between components. Operator strength can thus be dominated by encoding self-energy. These norms are not performance bounds. Explicit pair operators also require entangling operations whose costs cannot inherit the added local layer's absence of extra two-qubit gates. See S, branch-1 operator note and joint carry results.

\section{Readout, forces and ensemble uncertainty}

For the equal-field quadratic $U$, a fixed distribution of dipoles with mean $\bar\mu$ and covariance $\Sigma$ satisfies
\begin{equation}
\mathbb E U(\mu)=U(\bar\mu)+\frac12\operatorname{tr}(A\Sigma).
\end{equation}
For the mean $\hat\mu_M$ of $M$ independent shots from that fixed distribution,
\begin{equation}
\mathbb E U(\hat\mu_M)=U(\bar\mu)+\frac{\operatorname{tr}(A\Sigma)}{2M}.
\end{equation}
These displayed formulas concern $U$, not the mixed-field estimator $L$. The corresponding scaled contract must retain the declared normalisation of $W_c$. Optimising the energy of a decoded mean changes the objective from the usual expected QUBO cost; the mean can lie off the grid. The output rule must distinguish a selected bitstring, an expected cost and a sample-mean vector.

In the scalar example at $t=0.51$, the mean equals the continuous optimum, but finite-shot mean-energy biases are $0.62475/M$ for independent bits and $0.12495/M$ for the correlated mixture. The retained script enumerates independent-bit outcomes for $M=1,2,3$. These are arithmetic checks, not hardware sampling measurements. Training randomness, adaptive selection and correlated noise require separate analysis. See S, product-control note and joint investigation.

For mixed-field readout, let
\begin{equation}
C_{\mu z}=\mathbb E[(\mu-\bar\mu)(z-\bar z)^T].
\end{equation}
The supplied mixed-field note gives
\begin{equation}
\mathbb E L(\mu,z)=L(\bar\mu,\bar z)+\frac12\operatorname{tr}(AC_{\mu z}).
\end{equation}
Paired primal and adjoint samples can therefore contribute cross-covariance. Independent batches remove this contribution in expectation; they do not make each realised corrected energy accurate. Shared or adaptively coupled batches require covariance accounting. Residual certification of the realised vectors remains a separate check, with its own matrix and spectral-bound requirements. This identity is a conditional derivation, not a supplied shot experiment or a complete sampling analysis for a trained workflow.

Energy accuracy alone does not certify forces. For a smooth coordinate $R$, the corrected readout satisfies
\begin{align}
\frac{dL}{dR}&=\frac12\left[z^TA_R\mu-(E_0)_R^T\mu-(E_s)_R^Tz\right]\\
&\quad+\frac12r_0^T\mu_R+\frac12r_s^Tz_R.
\end{align}
Both response terms vanish at exact stationarity; approximate vectors require their control. An exact minimum envelope over a fixed grid of continuous variational energies is continuous but may have force jumps. That alone does not imply a nonconservative force. Reusing the stationary half-interaction formula can instead create energy jumps. In the scalar mixed-field switch with $A=1$, $E_s=R$, $E_0=1$, primal levels zero and one, and fixed adjoint $z=0.9$, the corrected energy jumps by $-0.05$ at $R=1/2$. Fixed laboratory grids can also introduce rotational artefacts; molecular-frame grids remain an alternative requiring frame definitions and derivatives. See S, mixed-field note and branch observable audits.

For a deterministic approximate potential with a domain-wide bound $|\delta U|\leq\eta$, the canonical free-energy error obeys $|\delta F|\leq\eta$. If both hydration endpoints change, their absolute error bounds add. Nonnegative solution variational errors imply an upward hydration shift only with the vacuum Hamiltonian preserved. A snapshot certificate supplies neither a domain-wide guarantee nor sampling overlap, equilibration or uncertainty estimates. Quantum shots are not the canonical molecular ensemble.

\section{Smallest useful experiment, resources and falsifier}

Acquire a publication-matched text fixture containing ordered $A$, $E_s$, $E_0$, reference dipoles and induction energy, with units, dielectric setting and exact implementation revision. Model bytes, solute and MM coordinates, numerical charges and cutoff/image rules are additionally needed for independent physical reconstruction. Displaced fixtures or derivatives are needed for forces. Historical access audits in S found concrete transport and compatibility blockers; those failures do not establish that advertised files are absent remotely.

First reconcile the reference energy and verify symmetry and positive definiteness. Select the equal-field, uniformly scaled or mixed-field contract. For each solver, record its training objective, decoded output rule and reported observable separately. A selected-vector benchmark can directly test the residual certificates above. If sample means or expected energies are also studied, treat them as explicitly labelled readout variants rather than silently substituting them for selected vectors.

For genuinely mixed fields, compare raw and corrected readouts from the same primal approximation against $H_*$. Declare how the adjoint is obtained and include all additional work. If both solves are quantum, optimise the two SPD quadratic objectives separately and evaluate $L$ afterwards. If the adjoint is classical, identify that hybrid allocation and compare against the classical primal solve with the same available factors. These are prospective benchmark choices; no completed implementation or measured resource saving is supplied.

Fix physical energy and force tolerances, grids, stopping rules and resource budgets before comparing the classical controls, both P circuit families and any separately justified pair-operator extension. Retain refinement, clipping, rigid rotations, grid-switch paths and finite-difference forces. A direct $3Nb$ encoding does not establish a useful bit precision or hardware capacity. The synthetic checks supply neither.

For an explicitly stored dense matrix of dimension $d$, a new Cholesky factorisation costs $O(d^3)$ and subsequent right-hand sides cost $O(d^2)$. Residual evaluation needs a matrix-vector product, not another solve. A Gershgorin pass and a residual evaluation each cost $O(d^2)$, with $O(d)$ norm work. A positive certified lower bound can support the residual certificate; a nonpositive Gershgorin result is inconclusive. Matrix assembly, spectral-bound construction and floating-point error control must be charged. An exact classical adjoint is not a free certificate: with factors available, the original primal solve also costs a further right-hand side. These inherited cost corrections are explicit in S, branch-3 ledger entries 13, 20 and 24.

Compare complete budgets: coefficient preparation, precision selection, initialisation, restarts, primal and adjoint optimisation, gradients and quantum metrics, shots, compilation, readout, certification and force recovery. Report success distributions and uncertainty rather than iteration counts alone. Keep model, discretisation, optimisation and statistical error separate.

The practical proposal fails if an apparent gain disappears against component-level classical updates, fails the required energy or force tolerance, or exceeds the original classical solve cost without compensating scientific value. A trainability benefit remains conjectural. Only after the snapshot and force checks pass should a defined molecular ensemble be sampled or reweighted with overlap and convergence assessment. The present evidence establishes neither useful quantum replacement nor molecular hydration improvement nor publication novelty. The pause means that the physical contract and accuracy-versus-cost comparison must be established before production integration; it is not a general impossibility claim.

\appendix
\section{Evidence and stable references}

\paragraph{Source boundary.}
Q is supplied as \texttt{sources/external/Q.txt}; its cover identifies a peer-reviewed repository manuscript and gives 5 September 2026 as the download date. P is supplied as \texttt{sources/external/P.txt}; its opening gives publication on 19 March 2024 and correction on 17 June 2024. The correction notice concerns interchanged Figures 3 and 5, and 4 and 6. The approved supplement's metadata identify P's underlying publisher PDF, but no PDF bytes are available in this transport. Q's primary PDF is also unavailable. Historical extraction checks are attributed records, not current visual inspection.

S denotes \texttt{sources/external/research-supplement.md}, an approved generated campaign artefact. Paths below are exact member keys in its \texttt{files} mapping. Its coverage contract explicitly limits completeness: scientific dependencies and some ledger content are omitted. Visibility, hashes and prior checks do not independently certify the scientific claims.

\paragraph{Inherited energy and circuit findings.}
Equal-field energy, residual, symmetry and endpoint arguments are in \texttt{branches/branch-1/ada/derivations.md}, \texttt{branches/branch-2/ada/observable-audit.md} and \texttt{branches/branch-3/ada/observable-audit.md}. Supplied stable ledger references include branch-2 \texttt{ledger.jsonl}, entries 7, 12 and 16--19, and branch-3 \texttt{ledger.jsonl}, entries 5, 9 and 13--15.

Circuit reachability and product controls are in \texttt{branches/branch-1/emmy/product-control-note.md} and \texttt{branches/branch-2/emmy/ansatz-results.json}, with branch-2 ledger entries 3, 5--6 and 16. Operator conventions and commutator checks are in \texttt{branches/branch-1/emmy/commutator-note.md} and \texttt{branches/branch-1/emmy/commutator-results.json}. The branch-1 ledger is cited within supplied notes but is not itself supplied here; those references are attributed pointers, not claims of current ledger inspection.

\paragraph{Implementation and acquisition uncertainty.}
The inspection report is \texttt{branches/branch-2/ada/v1-implementation-audit.md}, with branch-2 ledger entries 25--31. It concerns engine revision \texttt{787617c67391592466491efe276cc2ce267f0019}; historical equivalence to Q is unresolved. Scaling checks are \texttt{branches/branch-2/ada/v1\_scaling.py} and \texttt{branches/branch-2/ada/v1-scaling-results.json}. Access records are \texttt{branches/branch-1/ada/input-audit.md} and \texttt{branches/branch-2/emmy/v1-access-audit.json}. The inspected supporting-data revision is \texttt{0279424ec1770e155533af6247f28ebca98b72b5}. Script reading and failed binary transport did not supply a matched model and saved configuration.

\paragraph{Two-site evidence and corrected costs.}
Use \texttt{branches/branch-3/ada/two-site-audit.md}, \texttt{branches/branch-3/ada/two\_site.json} and \texttt{branches/branch-3/emmy/two\_site\_audit.json}, with branch-3 ledger entries 21--27. Preserve full enumeration versus pair-factorisation distinctions, trimmed rotated domains, refinement plateaux and the separate-field rounding failure. Entries 13, 20 and 24 correct the earlier claim that residual evaluation required a solve. The referenced Emmy encoding note and some original executables are omitted from this bounded transport; retained outputs and supplied notes must not be described as a complete runnable archive.

\paragraph{Joint results and reproduction limits.}
Mixed-field algebra and tests are \texttt{ada/mixed-field-note.md}, \texttt{ada/mixed\_field.py} and \texttt{ada/mixed\_field.json}. The script defines its rational matrix, prescribed vectors and fields internally and uses standard-library arithmetic. Its retained output records successful identity and bound assertions and a finite-difference tolerance check. Script hash prefix \texttt{969bc6a634ac} identifies the retained version. Joint \texttt{ledger.jsonl} entries 5 and 8 supply the result and clarification; entry 8 supersedes entry 5's notation and scope statement. \texttt{emmy/ada-mixed-field-check.json} records the partner's historical check.

Product-angle, carry-path and continuation results are \texttt{emmy/joint-investigation.md}, \texttt{emmy/carry\_control.py} and \texttt{emmy/carry-results.json}, with joint ledger entries 6, 9 and 12. The script defines its rational scalar and two-site inputs internally. Its retained hash prefix is \texttt{7a68c3c1e151}; it writes the retained result path when run from the research root. \texttt{ada/partner-check.json} records historical agreement. These execution descriptions document the frozen procedures and retained results, rather than fresh numerical evidence.

\texttt{ada/inherited-checks.json} contains ten historical replay records reporting passed assertions and no differences. \texttt{emmy/replay.py} records the script/output mappings and its comparison exclusions for hashes and dependency probes. Some inputs needed by those replays are omitted here. Historical replay of synthetic calculations does not establish molecular applicability, independent paper replication or fresh execution.

\paragraph{Prior art and unresolved scientific coverage.}
The supplied research attributes product-state optimisation to Bowles and colleagues, \url{https://arxiv.org/abs/2108.08064}; related Y-rotation and classical comparisons to Vijendran and colleagues, \url{https://arxiv.org/abs/2302.04479}; gauge-potential methodology to Claeys and colleagues, \url{https://arxiv.org/abs/1904.03209}; and adjoint correction to Venditti and Darmofal. Branch notes also identify established residual-force and truncated conjugate-gradient work. These are inherited, targeted attributions with recorded access limitations, not a new exhaustive literature review or priority determination.

The supplement records prior exposure to both P and Q. Neither is treated as unseen. Q's extracted Eq. 10 has an apparent repeated summation, and P has problematic extracted mathematical signs and inequalities. This account defines its conditional algebra independently rather than repairing those glyphs. Publication-matched inputs, historical implementation compatibility, molecular spectra, useful force tolerances, ensemble behaviour, trained quantum costs and novelty remain unresolved.

\end{document}
